\documentclass[11pt,a4paper]{article}

\usepackage[T1]{fontenc}
\usepackage[utf8]{inputenc}
\usepackage{amsmath,amssymb,amsthm}
\usepackage{booktabs}
\usepackage[margin=28mm]{geometry}
\usepackage[hidelinks]{hyperref}

\newtheorem{theorem}{Theorem}[section]
\newtheorem{lemma}[theorem]{Lemma}
\newtheorem{corollary}[theorem]{Corollary}
\newtheorem{proposition}[theorem]{Proposition}
\theoremstyle{definition}
\newtheorem{definition}[theorem]{Definition}
\newtheorem{remark}[theorem]{Remark}

\newcommand{\R}{\mathbb{R}}
\newcommand{\binomFtwo}{\binom{F}{2}}
\newcommand{\binomFthree}{\binom{F}{3}}

\title{Finiteness and Obstruction Bounds for\\
Face-Neighbourly Polyhedral Surfaces}
\author{Totocaedres project --- mathematical research note}
\date{21 November 2026}

\begin{document}
\maketitle

\begin{abstract}
We study closed orientable polyhedral surfaces embedded in $\R^3$ whose
planar faces are pairwise adjacent along at least one non-degenerate segment.
Faces may be non-convex, may have polygonal holes, and may meet the same face
along several disjoint segments; additional point contacts are allowed.  We
derive explicit necessary bounds that turn every fixed-genus search into a
finite combinatorial problem.  In particular, a genus-$g$ candidate has at
most $2g$ holes in all of its faces, and its number of faces is bounded by an
explicit quadratic inequality.

Every reduced vertex is an isolated intersection of three incident support
planes. This gives bounds on holes, contacts, and genus at a given face count,
including when several planes share a line. We prove that global finiteness
of the combinatorial search is equivalent to a universal bound on the number
of faces, or equivalently on the genus. We also give necessary conditions on
common support lines and on faces that straddle other support planes. The
existence of a universal bound remains unresolved. The geometric bounds
exclude exactly five faces even when holes and extra contacts are allowed.
\end{abstract}

\section{Model and notation}

Let $S\subset\R^3$ be the boundary of a compact connected set with non-empty
interior.  We assume that $S$ is a connected, closed, orientable, embedded
polyhedral $2$-manifold.  It is decomposed into $F\geq4$ connected planar
faces.  Face $i$ is a compact planar region of genus zero with
$b_i=h_i+1$ pairwise disjoint simple polygonal boundary components.  Every
boundary component has at least three edges.  Put
\[
    H=\sum_{i=1}^{F} h_i.
\]

The subdivision has $V$ vertices and $E$ edges.  Every edge belongs to two
distinct faces with different support planes. We count maximal planar faces
and exclude artificial seams between coplanar regions. Artificial degree-two
vertices have been suppressed, and a
face occurs at most once in the cyclic link of a vertex.  Consequently
\[
    3\leq \deg(v)\leq F.
\]
The principal problem considered here requires every pair of distinct faces
to share at least one positive-length edge.  For a pair $i<j$, let
\begin{itemize}
\item $m_{ij}$ be the number of subdivision edges incident with faces $i,j$;
\item $s_{ij}$ be the number of connected components of their shared-edge
      subgraph;
\item $p_{ij}$ be the number of vertices at which $i,j$ meet but no
      $ij$-edge is incident.
\end{itemize}
Thus $m_{ij}\geq s_{ij}\geq1$, while $p_{ij}$ may be positive in addition.
Define
\[
  M=E-\binomFtwo=\sum_{i<j}(m_{ij}-1),
  \qquad
  D=\sum_v(\deg(v)-3).
\]
Both $M$ and $D$ are non-negative.  The genus of $S$ is denoted by $g$.

\section{Universal combinatorial bounds}

The first results do not use coordinates, supporting planes, convexity, or a
general-position assumption.

\begin{theorem}[master identity]\label{thm:master}
Every surface in the model satisfies
\begin{equation}\label{eq:master}
  12g=(F-3)(F-4)+2M+2D+6H.
\end{equation}
\end{theorem}

\begin{proof}
A planar face with $h_i$ holes has Euler characteristic $1-h_i$.
Additivity of Euler characteristic for the face decomposition gives
\[
  V-E+F-H=2-2g.
\]
The degree sum is $2E=3V+D$.  Eliminating $V$, and then substituting
$E=\binomFtwo+M$, gives \eqref{eq:master}.
\end{proof}

\begin{theorem}[finite fixed-genus layer]\label{thm:fixed-genus}
Fix $g\geq0$ and define
\[
  R(g,F)=12g-(F-3)(F-4).
\]
Every feasible signature satisfies
\begin{align}
 F &\leq
 \left\lfloor\frac{7+\sqrt{48g+1}}{2}\right\rfloor,
 \label{eq:Fbound}\\
 R(g,F)&\geq0,\label{eq:Rpositive}\\
 M+D+3H&=\frac{R(g,F)}2,\label{eq:budget}\\
 H&\leq\left\lfloor\frac{R(g,F)}6\right\rfloor\leq2g,
 \label{eq:Hgenus}\\
 M,D&\leq\frac{R(g,F)}2.\label{eq:MDgenus}
\end{align}
In particular, only finitely many integer tuples $(F,M,D,H,E,V)$ occur at a
fixed genus.
\end{theorem}

\begin{proof}
All terms following $(F-3)(F-4)$ in \eqref{eq:master} are non-negative.
Therefore $(F-3)(F-4)\leq12g$, whose upper root is
\eqref{eq:Fbound}.  The remaining statements follow by rearranging
\eqref{eq:master}.  Finally,
\[
  E=\binomFtwo+M,
  \qquad
  V=\frac{2E-D}{3},
\]
so $E$ and $V$ also range over finite sets.
\end{proof}

\begin{corollary}[hole obstruction]\label{cor:holes}
No genus-$g$ candidate in the principal problem can have more than $2g$
polygonal holes in all faces combined.  More sharply, after $F$ is fixed,
\[
 H\leq
 \left\lfloor
 2g-\frac{(F-3)(F-4)}6
 \right\rfloor.
\]
Every unit of $M$ or $D$ consumes the same budget as one unit on the
right-hand side of \eqref{eq:budget}, while one face hole consumes three
units.
\end{corollary}

This is already enough to make a complete search through all genera
$0\leq g\leq G$ finite.  It does not merely bound a list of numerical
parameters: for fixed $E$ there are finitely many involutions pairing the
$2E$ darts, finitely many assignments of darts to faces, and finitely many
boundary permutations.  Quotienting by relabelling reduces this finite set
but is not required for finiteness.

\subsection{Additional universal filters}

\begin{lemma}[contour capacity]\label{lem:contours}
Every candidate satisfies
\begin{equation}\label{eq:contourcapacity}
  3(F+H)\leq2E=F(F-1)+2M.
\end{equation}
Equivalently,
\[
 H\leq\left\lfloor\frac{2E}{3}\right\rfloor-F,
 \qquad
 M\geq
 \left\lceil\frac{3H-F(F-4)}2\right\rceil
\]
whenever the latter lower bound is positive.
\end{lemma}

\begin{proof}
There are $F+H=\sum_i b_i$ boundary components.  Each contains at least
three edge occurrences, and every edge occurs in exactly two face
boundaries.
\end{proof}

\begin{lemma}[degree capacity]\label{lem:degree}
One has
\[
  D\leq(F-3)V,
  \qquad
  V\geq\left\lceil\frac{2E}{F}\right\rceil,
  \qquad
  D\leq\left\lfloor\frac{2E(F-3)}{F}\right\rfloor.
\]
Moreover $2E-D$ must be a non-negative multiple of three.
\end{lemma}

\begin{proof}
No face repeats in a vertex link, so $\deg(v)\leq F$.  Summing this
inequality and using $2E=3V+D$ proves the claims.
\end{proof}

The inequalities in Lemmas~\ref{lem:contours} and~\ref{lem:degree} should be
applied before any rotation-system enumeration.  They reject necessary
integer signatures, not merely geometric guesses.

\section{Support-plane bounds with degeneracies}\label{sec:planes}

Let $\Pi_i$ be the supporting plane of face $i$. The contact requirement and
the convention on coplanar seams imply that every two support planes are
distinct and nonparallel. Several planes can share a line or meet at a
single point.

\begin{lemma}[finite vertex set]\label{lem:vertices}
Every vertex is the unique intersection of a triple of incident support
planes. In particular,
\begin{equation}\label{eq:Varr}
  V\leq\binomFthree.
\end{equation}
For a fixed face $i$, the number $N_i$ of vertices on all its boundary
components satisfies
\begin{equation}\label{eq:Niface}
  N_i\leq\binom{F-1}{2}.
\end{equation}
\end{lemma}

\begin{proof}
The normals of the planes incident with a vertex have rank at least two.
If their rank were two, all the incident planes would contain the same
line through the vertex. Every incident edge would then lie on that line:
its two incident planes are distinct and intersect in that line. At most two
edges with disjoint relative interiors can emanate from a point along one
line, one in each direction. This contradicts $\deg(v)\geq3$. Thus the
incident normals have rank three.

Assign to each vertex the lexicographically first independent triple of
incident planes. The triple determines a unique point, so the assignment
is injective. For a vertex on face $i$, extend the normal of $\Pi_i$ to a
basis using two of the other incident normals. The resulting pair identifies
the vertex uniquely on $\Pi_i$.
\end{proof}

\begin{theorem}[holes per face]\label{thm:arr-holes}
Put
\[
  q_F=\left\lfloor\frac{(F-1)(F-2)}6\right\rfloor.
\]
Every face and the total surface satisfy
\begin{equation}\label{eq:arr-holes}
  b_i\leq q_F,
  \qquad
  h_i\leq q_F-1,
  \qquad
  H\leq F(q_F-1).
\end{equation}
\end{theorem}

\begin{proof}
The boundary components of one face are disjoint simple polygonal cycles, so
their vertex sets are disjoint.  Every cycle has length at least three.
Consequently $3b_i\leq N_i$.  Apply \eqref{eq:Niface} and sum over the
faces.
\end{proof}

Theorem~\ref{thm:arr-holes} bounds the number of holes in terms of $F$ alone.

\begin{theorem}[capacity of a face pair]\label{thm:paircapacity}
For every pair $i<j$, let $L_{ij}=\Pi_i\cap\Pi_j$ and let $k_L$ denote the
number of support planes containing a given intersection line $L$. Then
\begin{align}
  m_{ij}&=s_{ij},\label{eq:edgepair}\\
  2m_{ij}+p_{ij}&\leq F-k_{L_{ij}}\leq F-2.\label{eq:contactcapacity}
\end{align}
\end{theorem}

\begin{proof}
At a reduced vertex the cyclic link is a cycle of length at least three of
distinct face labels. An unordered pair of labels can occur at most once as
consecutive labels, so two edges shared by the same pair cannot share an
endpoint. This proves $m_{ij}=s_{ij}$. Their $2m_{ij}$ endpoints and the
$p_{ij}$ isolated contact vertices are all distinct and lie on $L_{ij}$.
The rank-three argument shows that each such vertex lies on at least one
incident plane that does not contain the line. Each of the $F-k_{L_{ij}}$
planes outside the pencil contributes at most one point on that line.
\end{proof}

\begin{corollary}[capacity of a common line]\label{cor:pencil}
For any intersection line $L$ contained in $k_L$ support planes,
\[
  \sum_{\{i,j\}:L_{ij}=L}m_{ij}\ \leq\ F-k_L-1,
  \qquad
  F\ \geq\ k_L+\binom{k_L}{2}+1.
\]
\end{corollary}

\begin{proof}
Every edge on $L$ has endpoints among the at most $F-k_L$ intersection
points with non-pencil planes. The edge interiors are disjoint and every
edge uses at least one of the intervals between consecutive points. Each
of the $\binom{k_L}{2}$ pairs in the pencil requires its own edge.
\end{proof}

\begin{theorem}[fixed-face genus bound]\label{thm:fixed-face}
Put $C=\binom{F}{3}$ and define the integer bounds
\begin{align}
 E_*&=\min\left\{\binom{F}{2}\left\lfloor\frac{F-2}{2}\right\rfloor,
                    \left\lfloor\frac{3C}{2}\right\rfloor\right\},\\
 H_*&=\min\left\{F(q_F-1),\ \left\lfloor\frac{2E_*}{3}\right\rfloor-F\right\}.
\end{align}
Every realization satisfies $E\le E_*$, $H\le H_*$, and
\begin{equation}\label{eq:Gdisk}
 g\leq \left\lfloor 1-\frac{F}{2}+\frac{C}{4}+\frac{H}{2}\right\rfloor
 \leq \left\lfloor 1-\frac{F}{2}+\frac{C}{4}+\frac{H_*}{2}\right\rfloor.
\end{equation}
\end{theorem}

\begin{proof}
For each pair, the number of vertices incident with both faces is at most
$F-2$ by Theorem~\ref{thm:paircapacity}. Summing gives
\[
 \sum_v\binom{\deg(v)}{2}\leq \binom{F}{2}(F-2)=3C .
\]
For every integer $d\geq3$, $d\leq\binom{d}{2}$ and
$d/2-1\leq\binom{d}{2}/6$, the latter since $(d-3)(d-4)\geq0$.
The degree sum gives $2E\leq3C$ and $E-V\leq C/2$. The first bound in
$E_*$ follows by summing \eqref{eq:contactcapacity}. The contour capacity
\eqref{eq:contourcapacity} and Theorem~\ref{thm:arr-holes} bound $H$.
The genus estimate follows from $g=1+(E-V-F+H)/2$.
\end{proof}

If all vertices are trivalent, $p_{ij}=0$ for every pair, since every pair of
labels in a cyclic link of length three is consecutive.

\begin{table}[ht]
\centering
\small
\begin{tabular}{rrrrrrr}
\toprule
$F$ & $q_F$ & $H_*$ & $\lfloor(F-2)/2\rfloor$ & $E_*$ &
Genus bound & With $H=0$\\
\midrule
4  & 1  & 0   & 1 & 6   & 0   & 0\\
7  & 5  & 21  & 2 & 42  & 16  & 6\\
8  & 7  & 48  & 3 & 84  & 35  & 11\\
9  & 9  & 63  & 3 & 108 & 49  & 17\\
12 & 18 & 204 & 5 & 330 & 152 & 50\\
\bottomrule
\end{tabular}
\caption{Upper bounds with arbitrary support-plane degeneracies.}
\label{tab:bounds}
\end{table}

\begin{corollary}[Exclusion of five faces]\label{cor:five}
There is no realization with $F=5$ in the stated model.
\end{corollary}

\begin{proof}
Theorem~\ref{thm:paircapacity} gives $m_{ij}\leq1$, so $E=10$.
By Corollary~\ref{cor:pencil}, a line contained in three support planes
would require at least seven faces. Therefore any three of the five
planes have at most one point in common.

A vertex of degree five would make all five planes concurrent.
Lemma~\ref{lem:vertices} would then force every vertex to be that same
point, which is impossible. Two vertices of degree four would have at
least three incident face labels in common. They would give two common
points of the corresponding three planes, also impossible. All vertices
thus have degree three except for at most one of degree four, so $D\leq1$.
But $3V+D=2E=20$ requires $D\equiv2\pmod3$, a contradiction.
\end{proof}

\section{The global question}\label{sec:global}

\begin{theorem}[Equivalent forms of global finiteness]\label{thm:global}
In the stated model the following properties are equivalent:
\begin{enumerate}
\item The set of realizable combinatorial types is finite.
\item There is an absolute upper bound on $F$.
\item There is an absolute upper bound on $g$.
\end{enumerate}
\end{theorem}

\begin{proof}
A finite set of types supplies both bounds. An upper bound on $F$ bounds
$E$ by Theorem~\ref{thm:fixed-face}. A map with at most that many edges is
described by finitely many pairings and permutations of its darts and
groupings of contour cycles into faces. An upper bound on $g$ bounds $F$
by \eqref{eq:Fbound}. This proves the equivalence.
\end{proof}

Both the genus upper bound in \eqref{eq:Gdisk} and the Euler lower bound
$(F-3)(F-4)/12$ grow with $F$. Their comparison leaves arbitrarily large
values of $F$ compatible with the inequalities. The five-face exclusion
and the seven-face Szilassi example illustrate why an empty layer cannot
serve as a stopping criterion.

\subsection{A necessary number of straddling pairs}

Say that face $i$ \emph{straddles} plane $\Pi_j$ if it has points on both
open sides of that plane. Call a pair straddling if either face straddles
the plane of the other. Let $B$ count such pairs and let $T$ count triples
of support planes with a common line.

\begin{proposition}\label{prop:straddling}
For $F\geq5$,
\[
 B \ \geq\ \max\left\{0,\
 \left\lceil\frac{\binom{F}{2}}{10}-\frac{3T}{F-2}\right\rceil\right\},
 \qquad
 T\leq\frac{k_{\max}-2}{3}\binom{F}{2},
 \quad
 k_{\max}=\max_L k_L .
\]
In particular, $B\geq F^2/20-O(F^{3/2})$ as $F$ increases.
\end{proposition}

\begin{proof}
Consider five faces with no straddling pair and no triple of support
planes with a common line. Relative to the plane of any one of the
five, each of the other four is entirely in a closed halfspace. If two
were in opposite halfspaces, their shared segment would lie in three
support planes, contrary to the choice of the five. Consequently each
of the five planes supports the convex hull of the five faces. This
hull is three-dimensional because the faces have positive area and
distinct planes. Each face lies in a separate hull facet. Every shared
segment lies on an edge between the corresponding facets. This would
give a $K_5$ in the planar dual graph of a convex polyhedron.

Every five-element set must therefore contain a straddling pair or a
triple with a common line. Counting the five-element sets containing
each of these yields
\[
 \binom{F}{5}\leq B\binom{F-2}{3}+T\binom{F-3}{2}.
\]
Each pair of planes determines exactly one line, so
\[
 T=\sum_L \binom{k_L}{3}
   \leq \frac{k_{\max}-2}{3}\sum_L\binom{k_L}{2}
   =\frac{k_{\max}-2}{3}\binom{F}{2}.
\]
Corollary~\ref{cor:pencil} gives $k_{\max}=O(\sqrt F)$.
\end{proof}

This extends the halfspace argument used for convex polygons in
\cite{ArsenevaEtAl2024}. It gives a necessary condition for any unbounded
family. A global upper bound by this route would require an incompatible
upper estimate on straddling pairs from enclosure and connectedness.

\subsection{An exact formulation using bounded chambers}

\begin{proposition}\label{prop:chambers}
The bounded solid enclosed by any realization is a union of closures of
bounded three-dimensional chambers of its supporting-plane arrangement.
\end{proposition}

\begin{proof}
A chamber is a connected component of the complement of the union of
the planes. As the boundary of the solid lies in the union of the planes,
each chamber lies entirely in its interior or entirely outside. An
interior chamber is bounded. The interior of the solid is open and the
union of the planes has empty interior. Interior chamber points are
therefore dense in the solid, which is the closure of its bounded
interior. There are finitely many chambers, so taking closures gives the
claimed union.
\end{proof}

For a fixed arrangement one can enumerate chamber selections, take the
boundary, and check the vertex links, connected planar regions and shared
edges. A search across all $F$ would require a theorem uniform in the
number of planes. The construction of arbitrary complete graphs as
side-contact representations in \cite{ArsenevaEtAl2024} uses polygonal
systems with boundary. Closing them while preserving complete adjacency
is an additional requirement.

\section{Consequences for exhaustive search}

\subsection{A provably finite fixed-genus scheduler}

For a requested maximum genus $G$, the following enumeration is complete at
the level of abstract maps.
\begin{enumerate}
\item For $0\leq g\leq G$, enumerate $4\leq F\leq F_{\max}(g)$ using
      \eqref{eq:Fbound}.
\item Compute $R(g,F)$ and enumerate the finite solutions of
      $M+D+3H=R/2$.
\item Reject signatures violating contour capacity
      \eqref{eq:contourcapacity}, degree capacity, or the integrality of
      $V=(2E-D)/3$.
\item Enumerate distributions of $M$ among the $\binomFtwo$ face pairs and
      distributions of $H$ among the $F$ faces.
\item Enumerate reduced rotation systems and quotient by face relabelling and
      global orientation reversal.
\end{enumerate}
All loops are finite for the chosen $G$.

\subsection{A provably finite fixed-face scheduler}

With fixed $F$, a run needs no input genus bound.
Theorems~\ref{thm:arr-holes}, \ref{thm:paircapacity}, and
\ref{thm:fixed-face} provide finite ranges for $H$, every $m_{ij}$, $E$, and
$g$. A finite face cap $F\leq F_*$ therefore gives a finite total search.

\subsection{Machine-checkable pruning rules}

The following filters are suitable for direct implementation before map
generation:
\[
\begin{array}{ll}
\text{universal:}&
(F-3)(F-4)\leq12g,\quad
M+D+3H=R/2,\\[2mm]
&3(F+H)\leq F(F-1)+2M,\quad
2E-D\equiv0\pmod3,\\[2mm]
&0\leq D\leq 2E(F-3)/F;\\[2mm]
\text{geometric:}&
h_i\leq q_F-1,\quad
2s_{ij}+p_{ij}\leq F-2,\\[2mm]
&m_{ij}=s_{ij},\quad
V\leq\binomFthree,\quad E\leq E_*.
\end{array}
\]
Every rejection can be accompanied by the violated inequality, producing a
small auditable impossibility certificate.

\section{Geometric decision and stronger obstructions}

For a fixed finite map, geometric realizability can in principle be written
in the first-order theory of the real numbers.  Use variables for every vertex
coordinate and every supporting plane.  Face incidences give polynomial
equalities; distinct vertices, non-zero edge lengths, polygon simplicity,
orientation, and the absence of unintended intersections give a finite
Boolean combination of polynomial strict and weak inequalities. A finite
disjunction over triangulations of the polygonal regions, including regions
with holes, lets us express intersection conditions with triangle
predicates. Quantifier elimination over real closed fields gives a
termination theorem for each fixed map.

This observation is theoretical rather than a practical algorithm: cylindrical
algebraic decomposition can be doubly exponential.  A realistic obstruction
pipeline should instead proceed as follows.
\begin{enumerate}
\item Apply the integer and arrangement bounds in this note.
\item Encode signs of $4\times4$ affine determinants as a rank-four oriented
      matroid feasibility problem.
\item Use SAT/SMT to reject incompatible sign systems.
\item Use exact nonlinear real arithmetic only for the surviving maps.
\item Archive either exact coordinates or a solver-verifiable unsatisfiable
      core.
\end{enumerate}
Schewe used the oriented-matroid/SAT strategy to prove non-realizability of
minimal twelve-vertex triangulations of genus six.  That result concerns the
vertex-neighbourly dual problem. Geometric realizability of the face-dual
maps requires a separate argument \cite{Schewe2010,RostVigh2026}.

\section{Status of the absolute bound}

Theorem~\ref{thm:global} reduces the question to a single conjectural
bound. A universal bound would yield a finite search. A family with
unbounded $F$ would imply unbounded genus and disprove global finiteness.
The plane and chamber arguments here establish neither outcome.
General embedded polyhedral surfaces can have arbitrarily high genus
\cite{Ziegler2005}. Their use here would require complete face adjacency.
The genus-three examples \cite{Mizhaev2026,RostVigh2026} use non-convex faces.

For trivalent disk faces with exactly one edge per face pair, the next
combinatorial layer after $F=7$ is
\[
  F=12,\qquad E=66,\qquad V=44,\qquad g=6.
\]
There are 59 corresponding neighbourly triangulations in the dual
combinatorial problem. R\"ost and V\'igh report that realizability of any
of the twelve-face dual maps remains open \cite{RostVigh2026}. This is one
specific layer without a resolution in the classical disk-face setting.

\begin{thebibliography}{9}

\bibitem{ArsenevaEtAl2024}
E.~Arseneva, L.~Kleist, B.~Klemz, M.~L\"offler, A.~Schulz,
B.~Vogtenhuber, and A.~Wolff,
\emph{Adjacency Graphs of Polyhedral Surfaces},
Discrete \& Computational Geometry 71 (2024), 1429--1455.
\href{https://doi.org/10.1007/s00454-023-00537-6}{doi:10.1007/s00454-023-00537-6}.

\bibitem{Mizhaev2026}
R.~Mizhaev,
\emph{Integer Realization of an Equivelar Octahedron of Genus 3},
arXiv:2609.17700, 2026.
\url{https://arxiv.org/abs/2609.17700}.

\bibitem{RostVigh2026}
G.~R\"ost and V.~V\'igh,
\emph{A Second Eight-Faced Polyhedron in Which Every Two Faces Share an Edge},
arXiv:2609.32998, 2026.
\url{https://arxiv.org/abs/2609.32998}.

\bibitem{Schewe2010}
L.~Schewe,
\emph{Nonrealizable Minimal Vertex Triangulations of Surfaces: Showing
Nonrealizability Using Oriented Matroids and Satisfiability Solvers},
Discrete \& Computational Geometry 43 (2010), 289--302.
\href{https://doi.org/10.1007/s00454-009-9222-y}{doi:10.1007/s00454-009-9222-y}.

\bibitem{Ziegler2005}
G.~M.~Ziegler,
\emph{Polyhedral Surfaces of High Genus},
arXiv:math/0412093, 2004.
\url{https://arxiv.org/abs/math/0412093}.

\end{thebibliography}

\end{document}
